% Abstract -- DATASET / BENCHMARK paper.
%
% Every number here is measured and reproducible from this repo; the source of each is in
% notes/paper/abstract-numbers.md. Nothing in this file is a placeholder: if a figure changes,
% change it here AND there.
%
% Deliberately absent, each for a stated reason (see abstract-numbers.md):
%   * the corpus size          -- the release filter was only just fixed, so the count moves
%   * any "first" claim        -- the gap is a conjunction of four conditions, and a reviewer
%                                 will take them apart one at a time
%   * the six-fold some-layers growth factor -- it is a LOWER bound, 18 of 30 unsolved
%   * the verification ceiling -- belongs in limitations, not here

\begin{abstract}
Crease-pattern-to-sequence is a natural testbed for multi-step spatial planning, but the
available resources do not support it: existing origami datasets provide crease patterns
without sequences, final states without intermediate steps, or sequences that were never
released. We contribute a synthesized corpus in which every sample is a (crease pattern, fold
sequence) pair \emph{by construction} --- folded forward from a square and unfolded to recover
the pattern --- so the ground truth requires no annotation and no trust, together with exact
forward engines and exhaustive solvers for two action-space tiers.

The \textbf{all-layers} tier restricts each fold to move the whole stack rigidly. We show this
buys three properties, of which the third is not usually named: no self-intersection, no layer
ordering, and \emph{no tearing is possible at all} --- a fold across a line moves no two
connected pieces relative to each other. The \textbf{some-layers} tier lets a fold move a
contiguous run of layers at either face of the stack, reinstating all three and making tearing
a question about the sheet's connectivity that is decidable only in the original coordinates.

Two measurements frame what the corpus is for. First, a \textbf{ceiling}: over an existing
366-pattern collection, 327 (89.3\%) are \emph{proven} to admit no simple-fold sequence, so a
pooled ``percent solved'' on such data is capped at 10.7\% and reports the task definition
rather than the model. Second, an \textbf{asymmetry that compounds}: generation is linear in
fold count while search is not. Median search cost rises $45.5\times$ and $91.9\times$ per two
folds in the all-layers tier and $247.5\times$ in the some-layers tier, whose cost relative to
all-layers grows from $4.5\times$ at three folds to $24.7\times$ at five; at six folds 18 of 30
some-layers instances are unsolved within a million simulated folds. Both solvers double as
verified ground truth and as the pure-search query baseline any learned method must be measured
against.
\end{abstract}
